\documentclass[10pt,a4paper]{amsart}
\usepackage[margin=19mm]{geometry}
\usepackage{amsmath,amssymb,mathtools}
\usepackage[scr=rsfs,cal=euler]{mathalfa}
\usepackage{xfrac}
\usepackage[unicode,hidelinks,bookmarks=false]{hyperref}
\newcommand{\ie}{i.e.\ }
\newcommand{\cf}{cf.\ }
\newcommand{\resp}{respectively\ }
\newcommand{\Subsection}{\subsection}
\providecommand{\nobreakdash}{\nobreak\hbox{-}}
\setlength{\emergencystretch}{2em}
\multlinegap0pt
\usepackage{mathtools}
\usepackage{enumerate}
\usepackage{bm}
\usepackage{amssymb}
\newtheorem{lemma}{Lemma}
\title{Covering a square by congruent squares}
\author{Gy\"orgy D\'osa, Zsolt L\'angi, Zsolt Tuza}
\date{Source: arXiv:2601.16535v3 (CC BY 4.0)}
\begin{document}
\maketitle

\noindent\emph{Source and licence.} Covering a square by congruent squares. Version arXiv:2601.16535v3 (CC BY 4.0), distributed by arXiv under the Creative Commons Attribution 4.0 International licence, http://creativecommons.org/licenses/by/4.0/, as recorded in the arXiv licence metadata for this exact version and shown on the abstract page https://arxiv.org/abs/2601.16535v3. Full licence notice: \texttt{licenses/arxiv-2601.16535-CC-BY-4.0.txt}.\par\smallskip

\noindent\textit{Editorial scope.} Counted unit: the article's lemma lem:mainfor3 (source0.tex lines 332--335). The article's complete original proof of it is delivered with it (source0.tex lines 337--346, one \texttt{proof} environment of 10 lines). No mathematics is written, repaired or re-derived: the statement and the proof are the article's own lines, in the article's own notation, and no retained line is substituted or re-flowed.  No further source-local context is needed: every cross-reference of every retained line resolves inside the two delivered spans.  Nothing was omitted from any retained span, so the package carries no omission record.  No retained line contains a citation, so the wrapper carries no bibliography and no bibliography entry.  No retained line contains a figure environment, an image-inclusion command or drawing code, so no image is delivered and the package declares no asset dependency.  Editorial formatting only, in addition: an AMS \texttt{amsart} preamble that restates the article's own declarations and macros used by the retained lines, namely the environments \texttt{lemma} on separate counters, together with the standard packages those lines need.  The retained environments print in this document as Lemma 1 in order of appearance, whereas the article prints the same items with the numbering of its own full document.  The complete native source of the article as distributed in the arXiv e-print for this exact version is bundled unchanged as \texttt{sources/193237/source0.tex}.
\begin{lemma}\label{lem:mainfor3}
Let $X$ be an $L$-shape maximally embedded in a unit square $S$. Assume that the legs of $X$ are of lengths $x$ and $l(x)$, where $x > 1$. Then $l(x) < 1$, the endpoint $p$ of $X$ lying on the longer leg is a vertex of $S$, and the other endpoint and the corner of $X$ are points of the two sides of $S$ not containing $p$.
Furthermore, $l(x)=x-x\sqrt{x^2-1}$.
\end{lemma}

\begin{proof}
The statement that $l(x) < 1$ is trivial. Let $q$ be the endpoint of $X$ on the leg of length $l(x)$, and $c$ be the corner of $X$.
Without loss of generality, we may assume that $p,q,c$ belong to the boundary of $S$. Since $x > 1$, it follows that $c$ is an interior point of a side $E_2$ of $S$. Thus, we also have that $q$ belongs to the interior of a side $E_1$ consecutive to $E_2$. Let $E_3$ be the other consecutive side of $E_2$. If $p$ is in the interior of the fourth side $E_4$, then rotating $X$ around $c$ we find an $L$-shape in $S$ whose shorter leg is longer than $l(x)$, a contradiction.

Assume that $p$ is in the interior of $E_3$. If, in addition, $p$ is not closer to $E_2$ than $q$ is, then we may rotate $X$ around $q$ and find an $L$-shape in $S$ with a longer short leg. 
%Biralo 5
%On the other hand, if $p$ is closer to $E_2$ than $q$ is, then we may rotate $X$ around %the midpoint of the segment $[p,q]$ and find a larger $L$-shape. 
On the other hand, if $p$ is closer to $E_2$ than $q$ is, then we may rotate $X$ around $p$ and find a larger $L$-shape. 
Thus, by the definition of $l(x)$, $p$ is the common point of $E_3$ and $E_4$, proving the first statement. Then the formula for $l(x)$ follows by the Pythagorean theorem and comparing the sides of two similar triangles.
\end{proof}

\end{document}
